\documentclass[10pt]{beamer}
\usepackage[T2A]{fontenc}
\usepackage[utf8]{inputenc}
\usepackage[english,russian]{babel}
\usepackage{amssymb,amsfonts,amsmath,mathtext}
\usepackage{cite,enumerate,float,indentfirst}
\usepackage{ragged2e}
\usepackage{comment}
\usecolortheme{whale}
\usepackage{multirow}
\setbeamertemplate{caption}[numbered]
\usepackage{caption}
\usepackage{fontawesome}
\usepackage{media9}
\usepackage{graphicx, animate}


\setbeamercolor{footline}{fg=blue}
\setbeamertemplate{footline}{
  \leavevmode%
  \hbox{%
  \begin{beamercolorbox}[wd=.93\paperwidth,ht=2.25ex,dp=1ex,center]{}%
   {\small МНШ лаб. 69 «Разработка теоретических основ и методов практической реализации ЛНС привязной высотной беспилотной платформы»}
    %МНШ лаб. 69, руководитель д.т.н., проф., Вишневский В.М.
  \end{beamercolorbox}%
    \begin{beamercolorbox}[wd=.07\paperwidth,ht=2.25ex,dp=1ex,right]{}%
  \insertframenumber{} / \inserttotalframenumber \hspace*{2ex}
  \end{beamercolorbox}}%
  \vskip0pt%
}

\begin{document}

%надо куда-то поместить на слайд. м.б. в виде картинки?
%Федеральное государственное бюджетное учреждение науки\\
%Институт проблем управления им. В.А. Трапезникова\\
%Российской академии наук (ИПУ РАН)

\title{Разработка теоретических основ и методов практической реализации системы локальной навигации привязной высотной беспилотной платформы} 

\author{{\Large \bf ОТЧЕТ} \\
по работе молодежной научной школы лаб. 69  \\
за первый (2024-2025) год \\
\bigskip
Руководитель МНШ: д.т.н., проф., Вишневский Владимир Миронович
}

\date{Москва, 2025}
\logo{\includegraphics[height=0.8cm]{logo}}

\frame{\titlepage} 




\begin{frame}{Состав участников}


\begin{minipage}{0.4\textwidth}
\justifying
\begin{enumerate}
\item Абрамян В.Л., м.н.с., аспирант МФТИ
\item Иванов М.Г., инженер, магистрант 1 курс МЭИ
\item Иванова Н.М., к.ф.-м.н., с.н.с.
\item Калмыков Н.С., н.с., аспирант ИПУ РАН
\item Лашин М.В., инженер, аспирант ИПУ РАН
\item Новочадова А.В. инженер, магистрант 2~курс МИИГАиК
 \end{enumerate}
\end{minipage}
\hfill
\begin{minipage}{0.58\textwidth}
\begin{figure}[h]
\begin{tabular}{l}
\includegraphics[width=0.99\linewidth]{1.png}
\end{tabular}
\end{figure}
\end{minipage}


\end{frame}


\begin{frame}{Актуальность и цель исследования}
\justifying
\small

Привязные высотные беспилотные платформы играют ключевую роль в обеспечении связи, мониторинга и ретрансляции сигналов в труднодоступных регионах. Их эффективная эксплуатация требует точного позиционирования и стабилизации, особенно в условиях недоступности сигналов спутниковой навигации, ветровых нагрузок и изменяющихся атмосферных условий.

\smallskip

{\bf Исследуемая проблема}: обеспечение надежности функционирования платформы в условиях недоступности сигналов спутниковой навигации.

\smallskip

{\bf Предлагаемое решение}: разработка локальной навигационной системы (ЛНС) привязной платформы, которая обеспечит
\begin{itemize}
\justifying
    \item Высокую точность позиционирования -- корректировку положения платформы,
    \item Стабилизацию платформы -- компенсацию колебаний, вызванных ветром, для беспрерывной связи.
    \item Автономность работы -- снижение зависимости от GPS/ГЛОНАСС, которые могут быть уязвимы к помехам.
    \item Безопасность -- предотвращение аварийных ситуаций, связанных с потерей управления.
\end{itemize}

\end{frame}


\begin{frame}{Полученные за отчетный год основные результаты}
 \begin{enumerate}
 \justifying
 \begin{small}
\item Проведен анализ мировых тенденций, перспективных технологий и математических методов проектирования навигационных систем для автономных и привязных БПЛА.
\item Разработаны архитектура и принципы построения системы локальной навигации привязной высотной платформы.
\item Разработаны технические и программные средства реализации радиолокационной системы навигации. 
\item Разработаны принципы построения и алгоритмы работы локальной системы навигации на базе лазерного лидара. 
\item Разработаны методы оценки надежности системы локальной навигации на основе построения дерева рисков. 
\item Разработаны аналитических моделей для исследования характеристик гибридной системы локальной навигации. 
\item Разработан макетный образец радиолокационной системы навигации. 
\item Разработан макетный образец системы локальной навигации на базе лазерного лидара. 
\end{small}
 \end{enumerate}
\end{frame}

\begin{frame}{Степень выполнения поставленных задач проекта}
\justifying

{\bf План работ} первого года выполнен {\bf полностью}. В 2024-2025 гг. опубликованы  \underline{$33$ научные работы}, включая $1$ монографию, $16$ статей в изданиях индексируемых в WoS/Scopus (в том числе $2$ статьи в журналах квартиля Q1); принято участие в $9$-ти международных и всероссийских конференциях; получено $2$ свидетельства о государственной регистрации программ. 

\bigskip

Монография \textbf{\textit{«Теория очередей и машинное обучение»}} объявлена лауреатом общенационального конкурса «Учебник года»  в номинации «Компьютерные науки», а также признана лучшей монографией в конкурсе ИПУ РАН в 2024 г. Статья \textbf{\textit{ «Investigation of the Fork–Join System with Markovian Arrival Process Arrivals and Phase-Type Service Time Distribution Using Machine Learning Methods» }}
признана лучшей статьей в конкурсе ИПУ РАН в 2024 г.

\bigskip

Принято к печати $7$ публикаций в журналах ВАК (RSСI) и WoS/Scopus.

\smallskip
В 2024 году защищена 1 кандидатская диссертация.

\end{frame}

\begin{frame}[allowframebreaks]{Участие в научных мероприятиях по тематике проекта}
 \justifying
\begin{footnotesize}
 \begin{enumerate}
 \justifying
\item XIV Всероссийское совещание по проблемам управления, 17-20 июня 2024, г. Москва.
\item  ХХ Всероссийская школа-конференция молодых ученых «Управление большими системами», 10-13 сентября 2024 г., г. Новочеркасск.
\item 27th International Conference on Distributed Computer and Communication Networks DCCN 2024, 23-27 September 2024, Moscow, Russia.
\item The Fourth International Workshop on Stochastic Modeling and Applied Research of Technology, SMARTY 2024, August 26-30, 2024, Karelian Research Center of RAS, Petrozavodsk.
\item 8th International Conference on Information, Control, and Communication Technologies (ICCT 2024), Оctober 1-5, 2024, Vladikavkaz, Russian Federation.
\item IEEE International Conference on Advanced Networks and Telecommunications Systems, 15–18 December 2024 // Assam, India // IIT Guwahati.
\item AnalytiX-2024 (международный семинар). Сентябрь 2024. Central University of Rajasthan, Rajasthan, India 
\item 23rd International Conference named after A.F. Terpugov Information Technologies and Mathematical Modelling. Октябрь 2024. Карши, Узбекистан.
\item 1-й Международная научная конференция «Школа теории массового обслуживания», г. Томск, ТГУ, 21-26 апреля 2025 г.
\end{enumerate}
\end{footnotesize}

\bigskip

Проведено 12 семинаров лаб. 69 по обсуждению плана выполнения работ МНШ и результатов решения поставленных задач.

\bigskip

В 2024 г. члены МНШ приняли активное участие в подготовке и проведении
\begin{small}
\begin{itemize}
\justifying
    \item 27-ой международной конференции Distributed Computer and Communication Networks: Control, Computation, Communications (DCCN), которая проводилась в РУДН,
    \item Международной конференции «Информационно-телекоммуникационные технологии и математическое моделирование высокотехнологичных систем», организованной РУДН и ИПУ РАН
    \item 8-ой Международной конференции «Информационные технологии и технические средства управления» (ICCT-2024), Владикавказ. 
\end{itemize}  
\end{small}
\end{frame}

\begin{frame}{План работ второго года выполнения проекта}
\begin{enumerate}
\justifying
\begin{small}
\item Разработка алгоритма эмуляции сигналов спутниковой навигации и их передачи по каналу связи земля-борт в систему управления БПЛА.  
\item Разработка оптической системы локальной навигации привязной высотной платформы с использованием методов машинного обучения.  
\item Анализ вариантов использования RFID-технологии для построения системы локальной навигации.  
\item Разработка методов оценки надежности системы локальной навигации на основе построения дерева рисков.  
\item Разработка аналитических моделей для исследования характеристик гибридной системы локальной навигации.  
\item Разработка программы и методики испытаний разработанных макетных образцов локальной навигационной системы в лабораторных и полевых испытаниях на базе комплекса «Альбатрос». 
\item Проведение испытаний, сравнительный анализ и рекомендации по выбору оптимального варианта построения локальной навигационной системы в зависимости от условий эксплуатации привязной высотной платформы, назначения, а также в условиях функционирования. 
\end{small}
\end{enumerate}
\end{frame}

\begin{frame}
\centering 
 {\bf Результаты, полученные участниками МНШ}
\end{frame}


    
\begin{frame}{План дальнейших исследований}
  
  
  
\end{frame}


\end{document}

